\chapter{Composición de rutas, registros y firmas}
\label{chap:composicion-especies-rev6}
\index{composición de rutas}
\index{composición de registros}
\index{mesón!expresión composicional}
\index{barión!expresión composicional}

Una especie compuesta no puede entrar en la ley de masas como un nombre al
que se asigna retrospectivamente una fila de cinco enteros. Antes de la firma
debe existir un objeto que pueda componerse. El atlas del
\cref{chap:atlas-genealogia-masas} proporciona el alfabeto finito; la
extensión superviviente y el registro del
\cref{chap:extension-ruta-caracter} proporcionan la memoria que puede
componerse. La cuestión de este capítulo es construir el estrato intermedio
sin consultar masas objetivo:
\[
\begin{gathered}
 \text{configuración composicional tipada}
 \longrightarrow
 \text{registros completos de sus constituyentes}\\
 \longrightarrow
 \text{registro compuesto}
 \longrightarrow
 \ZZ^5.
\end{gathered}
\]

El orden de las flechas es parte del resultado. Las expresiones se forman
antes de asignarles una realización física; los registros se componen antes
de extraer la firma; la firma se evalúa mediante el carácter sólo después de
haber compuesto la memoria. De este modo, el dominio compuesto no queda
reducido a las especies incluidas en una tabla metrológica.

\section{Clasificación interna y gramática composicional}
\label{sec:alfabeto-composicional-rev6}

Sea
\[
 \mathcal C_{81}=I_9\times I_9,
 \qquad I_9=\{1,\ldots,9\},
\]
el atlas de celdas, y sea
\[
 \rho(r,c)=(10-r,10-c)
\]
su inversión central. Denotamos por
\(\mathcal L,\mathcal N,\mathcal D,\mathcal U\) los subconjuntos de
leptones cargados, neutrinos, quarks de tipo abajo y quarks de tipo arriba,
respectivamente, y ponemos
\[
 \mathcal Q=\mathcal D\sqcup\mathcal U.
\]
Cada celda posee un nivel combinatorio
\(\operatorname{gen}(x)\). En el sector quark se conserva además el residuo
de color
\[
 \kappa(x)=r(x)-c(x)\pmod3.
\]
Escribimos
\[
 \mathcal Q_a=\{q\in\mathcal Q:\kappa(q)=a\},
 \qquad a\in\mathbb Z/3\mathbb Z.
\]
Por la \cref{prop:color-atlas},
\[
 |\mathcal Q_0|=|\mathcal Q_1|=|\mathcal Q_2|=12.
\]

\begin{definition}[Dual formal e identidad de celda]
\label{def:dual-identidad-composicional}
Para una sección de celda orientada \(x\), su dual formal \(x^\vee\) invierte
la orientación de carga y de color, y tiene soporte \(\rho(x)\). La
notación conserva la distinción
\[
 x^\vee\ne \rho(x)
\]
cuando \(\rho(x)\) se considera como celda sin orientación. Para cada
\(x\in\mathcal C_{81}\), denotamos por
\(\operatorname{id}_x:x\to x\) la flecha identidad. El conjunto de
identidades es
\[
 \mathfrak I
 :=\{\operatorname{id}_x:x\in\mathcal C_{81}\}.
\]
\end{definition}

La identidad \(\operatorname{id}_x\) es una flecha neutra del sistema
composicional.
No es, por esa sola condición, una realización de un bosón neutro. Del mismo
modo, el dual formal incorpora orientación y no puede sustituirse por una
segunda celda desnuda.

\begin{definition}[Expresiones mesónicas y bariónicas]
\label{def:expresiones-mesonicas-barionicas}
El conjunto de expresiones mesónicas neutras en color es
\[
 \mathfrak M
 :=
 \left\{
   M(q,q')=q\otimes{(q')}^\vee:
   q,q'\in\mathcal Q,\ \kappa(q)=\kappa(q')
 \right\}.
\]
Los pares son ordenados. Su residuo total es
\(\kappa(q)-\kappa(q')=0\).

El conjunto de expresiones bariónicas con posiciones de color fijadas es
\[
 \mathfrak B
 :=
 \mathcal Q_0\times\mathcal Q_1\times\mathcal Q_2,
\]
y escribimos
\[
 B(q_0,q_1,q_2)=q_0\otimes q_1\otimes q_2.
\]
Su neutralidad residual es la identidad
\[
 0+1+2=0\pmod3.
\]
\end{definition}

Estas expresiones son objetos sintácticos tipados. En particular, no se ha
tomado todavía un cociente por permutaciones, espín, sabor, estadística,
energía de enlace o equivalencia dinámica.

\begin{definition}[Flechas cargadas compatibles]
\label{def:flechas-cargadas-compatibles}
Las flechas leptónicas orientadas son
\[
 \mathfrak A_\ell
 :=
 \left\{
 (x,y)\in
   (\mathcal L\times\mathcal N)
   \sqcup
   (\mathcal N\times\mathcal L):
 \operatorname{gen}(x)=\operatorname{gen}(y)
 \right\}.
\]
Las flechas quark orientadas son
\[
 \mathfrak A_q
 :=
 \left\{
 (x,y)\in
   (\mathcal U\times\mathcal D)
   \sqcup
   (\mathcal D\times\mathcal U):
 \begin{array}{l}
 \operatorname{gen}(x)=\operatorname{gen}(y),\\[-2pt]
 \kappa(x)=\kappa(y)
 \end{array}
 \right\}.
\]
Ponemos
\[
 \mathfrak A_{\rm ch}:=\mathfrak A_\ell\sqcup\mathfrak A_q.
\]
\end{definition}

La igualdad de nivel tipa la transición en ambos sectores; la igualdad de
color se exige además en el sector quark. La orientación distingue \(x\to y\)
de \(y\to x\).

\begin{definition}[Sistema finito de composiciones tipadas]
\label{def:gramatica-celular-compuesta}
El sistema composicional finito es la unión disjunta tipada
\[
 \mathfrak G_{\rm comp}
 :=
 \mathcal C_{81}
 \sqcup\mathfrak I
 \sqcup\mathfrak M
 \sqcup\mathfrak B
 \sqcup\mathfrak A_{\rm ch}.
\]
Las \(81\) celdas son los generadores primitivos. Las \(81\) identidades son
flechas sobre esos mismos generadores y no un segundo atlas.
\end{definition}

\section{Censo combinatorio de composiciones de rutas independiente de datos metrológicos}
\label{sec:censo-composicional-rev6}

\begin{theorem}[Censo del sistema finito de composiciones]
\label{thm:censo-gramatica-composicional-rev6}
El sistema de la \cref{def:gramatica-celular-compuesta} contiene
\[
 \boxed{\begin{gathered}
        81\ \text{identidades},\qquad
        432\ \text{expresiones mesónicas},\\
        1728\ \text{expresiones bariónicas},\qquad
        372\ \text{flechas cargadas}
 \end{gathered}}
\]
Las flechas cargadas se descomponen exactamente como
\[
 372=320_{\text{leptónicas}}+52_{\text{quark}}.
\]
El censo depende sólo del atlas, de su inversión, de los niveles y del
residuo de color. No utiliza masas, unidades ni firmas publicadas.
Los cuatro cardinales cuentan expresiones composicionales internas o flechas
admisibles del sistema; no cuentan especies físicas distintas, partículas
físicas ni masas.
\end{theorem}

\begin{proof}
La asignación \(x\mapsto\operatorname{id}_x\) es biyectiva, de modo que
\[
 |\mathfrak I|=|\mathcal C_{81}|=81.
\]
Para los mesones, la condición de neutralidad permite escoger
independientemente un par ordenado dentro de cada clase de color. Como cada
clase contiene doce celdas,
\[
 |\mathfrak M|
 =\sum_{a\in\mathbb Z/3\mathbb Z}|\mathcal Q_a|^2
 =3\cdot12^2
 =432.
\]
Para los bariones, las posiciones \(0,1,2\) están fijadas, y por tanto
\[
 |\mathfrak B|
 =|\mathcal Q_0||\mathcal Q_1||\mathcal Q_2|
 =12^3
 =1728.
\]

Las familias leptónicas cargadas tienen tamaños
\[
 |\mathcal L_{G0}|=1,\qquad
 |\mathcal L_{G2}|=8,\qquad
 |\mathcal L_{G4}|=16,
\]
mientras las familias neutrínicas tienen tamaños
\[
 |\mathcal N_{G1}|=2,\quad
 |\mathcal N_{G2}|=4,\quad
 |\mathcal N_{G3}|=6,\quad
 |\mathcal N_{G4}|=8.
\]
Sólo \(G2\) y \(G4\) son niveles comunes. Como se cuentan ambas
orientaciones,
\[
 |\mathfrak A_\ell|
 =2(8\cdot4+16\cdot8)
 =320.
\]

Para las flechas quark, en el orden de colores \((0,1,2)\), las
multiplicidades de tipo arriba y tipo abajo en los niveles comunes son
\[
\begin{array}{ccc}
 &\mathcal U&\mathcal D\\
 G1&(2,1,1)&(0,1,1)\\
 G3&(4,4,4)&(2,2,2).
\end{array}
\]
El número de pares de una orientación y del mismo color es, por ello,
\[
 (2\cdot0+1\cdot1+1\cdot1)
 +(4\cdot2+4\cdot2+4\cdot2)
 =2+24=26.
\]
Al incluir la orientación contraria se obtiene
\[
 |\mathfrak A_q|=2\cdot26=52.
\]
Finalmente,
\[
 |\mathfrak A_{\rm ch}|=320+52=372.
\]
Ninguna de estas reglas contiene un valor objetivo o un vector de masa.
\end{proof}

\begin{corollary}[Constructor composicional sin blancos]
\label{cor:constructor-composicional-sin-blancos}
El constructor que enumera conjuntamente identidades, expresiones mesónicas,
expresiones bariónicas y flechas cargadas factoriza por los campos tipados
\[
 (\text{sector},\text{familia},\text{nivel},\kappa,\rho)
\]
del atlas. Con la notación anterior, sus cuatro codominios son
\(\mathfrak I,\mathfrak M,\mathfrak B,\mathfrak A_{\rm ch}\).
En particular, el censo puede fijarse antes de abrir cualquier tabla de firmas.
\end{corollary}

\begin{proof}
Las cuatro definiciones y el cálculo del
\cref{thm:censo-gramatica-composicional-rev6} emplean únicamente esos campos.
La tabla de firmas no pertenece al dominio de ninguna de las reglas.
\end{proof}

La independencia de blancos no significa que las expresiones hayan adquirido
ya un nombre físico. Significa algo anterior y más preciso: existe un dominio
composicional construido sin usar como criterio la cantidad que después se
evaluará.

\section{Transporte de memoria desde la ruta compuesta hasta la firma integral}
\label{sec:memoria-composicion-especies-rev6}

Sea \(\gamma\in\Gamma_{108}^{\rm enr}\) una presentación enriquecida de una
extensión. Su registro leído contiene
\[
 \mathcal E(\gamma)
 =
 (n_A,e_0,\ldots,e_{11},T_\zeta,C_\zeta),
\]
además de los datos de fase, hoja, orientación y borde que tipan una composición.
Los doce eventos no son reemplazables por su total antes de componer. Para
\(0\le i<6\), definimos
\[
 p_i=e_i+e_{i+6},
 \qquad
 a_i=e_i-e_{i+6},
\]
\[
 s_C=\sum_{i=0}^{5}p_i,
 \qquad
 M_i=p_i-p_5\quad(0\le i<5).
\]
Con \(M_5=(M_0,\ldots,M_4)\) y
\(A_6=(a_0,\ldots,a_5)\), la reconstrucción es
\[
 p_5=\frac{s_C-\sum_{i=0}^{4}M_i}{6},
 \qquad
 p_i=p_5+M_i,
\]
\[
 e_i=\frac{p_i+a_i}{2},
 \qquad
 e_{i+6}=\frac{p_i-a_i}{2}.
\]
Por el \cref{thm:memoria-1-5-6}, sobre la imagen integral se obtiene la
descomposición reversible
\[
 \boxed{(s_C,M_5,A_6),\qquad 1+5+6=12.}
\]
El primer bloque conserva el total visible; el segundo, la forma relativa
dentro de la hoja simétrica; el tercero, la memoria entre las dos semicoronas.

La coordenada torsional fuerte requiere además la cocadena a nivel de paso:
\[
 k_{\Delta_4}=T_\zeta-2C_\zeta.
\]
No es una función del mero patrón de eventos nulos. Esta diferencia será
esencial al interpretar el rombo de no descenso.

\begin{definition}[Dato de composición orientada]
\label{def:dato-empalme-orientado-rev6}
Sean \(\Gamma,\Lambda\) dos registros enriquecidos. Un dato de composición es un
par \((g,B)\), donde \(g\) alinea fase, hoja, orientación y posición en la
interfaz, y \(B\) es la subtraza común, con su contribución a la acción, sus
doce eventos, sus predecesores torsionales y su corona. Sobre el dominio en
que esos datos son compatibles se define
\[
 \boxed{
 {(n,e)}_\Gamma\star_{(g,B)}{(n,e)}_\Lambda
 =
 \bigl(
 n_\Gamma+n_\Lambda-n_B,\,
 e_\Gamma+g e_\Lambda-e_B
 \bigr).}
\]
La palabra torsional y la corona se componen en el mismo nivel, antes de
calcular \(T_\zeta,C_\zeta\) y \(k_{\Delta_4}\).
\end{definition}

La fórmula es una inclusión-exclusión de historias sobre una subtraza
identificada. Está demostrada con estructura de partida explícita: los
registros, el transporte \(g\), el borde \(B\), la orientación y la palabra
de predecesores forman parte de la entrada. El sistema composicional no elige por
sí sola ninguno de esos datos.

\begin{definition}[Extractor posterior a la composición]
\label{def:extractor-posterior-empalme-rev6}
Para un registro enriquecido completo se define
\[
 L_{\rm tick}(\gamma)
 =
 \bigl(
 n_A,s_C,k_{\Delta_4},
 \nu_{120}^{\rm ev},\nu_{270}
 \bigr)\in\ZZ^5,
\]
donde, si \(\tau_m=\operatorname{sgn}(e_m)\) e
\(I_a=\{a,a+4,a+8\}\pmod {12}\),
\[
 \nu_{120}^{\rm ev}
 =
 \sum_{a=0}^{3}
 \operatorname{sgn}\left(\sum_{m\in I_a}e_m\right),
 \qquad
 \nu_{270}
 =
 \sum_{m=0}^{11}\tau_m\tau_{m+3}.
\]
La firma de una composición se define por
\[
 \operatorname{Sig}(\Gamma,\Lambda;g,B)
 :=
 L_{\rm tick}
 \bigl(\Gamma\star_{(g,B)}\Lambda\bigr).
\]
\end{definition}

Los signos y las autocorrelaciones se aplican al registro ya compuesto. Por
ello, la definición no equivale a sumar las firmas de los constituyentes.

\section{Diagrama conmutativo de composición y preservación de la fibra}
\label{sec:rombo-no-descenso-composicional-rev6}

El catálogo histórico B1 emplea una proyección de eventos que conviene
tipar por separado. Sea
\[
 \Sigma_{12}=(+,+,+,-,-,-,+,+,+,-,-,-),
 \qquad \tau_m=\operatorname{sgn}(e_m),
\]
y definamos la proxy de eventos nulos
\[
 K_0(e)
 =
 \sum_{\tau_m=0}\Sigma_{12}(m).
\]
La proyección exacta es
\[
 F_0(n_A,e)
 =
 \left(
 n_A,\,
 \sum_m e_m,\,
 K_0(e),\,
 \sum_{a=0}^{3}
   \operatorname{sgn}\sum_{m\in I_a}e_m,\,
 \sum_m\tau_m\tau_{m+3}
 \right)\in\ZZ^5.
\]
Para dos registros alineados, escribimos
\[
 (n,e)\boxplus(n',e')=(n+n',e+e').
\]
Esta operación es el caso sin subtraza común y con alineamiento identidad del
caso de composición de eventos.

\begin{proposition}[Obstrucción al descenso de la proyección B1]
\label{prop:no-descenso-proyeccion-b1-rev6}
No existe una operación binaria
\[
 \mu:\ZZ^5\times\ZZ^5\longrightarrow\ZZ^5
\]
tal que
\[
 F_0(\Gamma\boxplus\Lambda)
 =
 \mu\bigl(F_0(\Gamma),F_0(\Lambda)\bigr)
\]
para todos los registros B1.
\end{proposition}

\begin{proof}
Consideremos las rutas, escritas en la carta
\((r,c,h_0,v_0)\),
\[
 a=(1,3,-1,+1),
 \qquad
 b=(6,7,-1,+1),
 \qquad
 c=(1,2,-1,-1).
\]
Las dos primeras tienen registros de eventos distintos:
\[
\begin{aligned}
 e(a)&=(5,4,4,-4,-4,-6,\;5,4,4,-4,-4,-6),\\
 e(b)&=(4,5,4,-4,-6,-4,\;4,5,4,-4,-6,-4),
\end{aligned}
\]
pero publican la misma cinco-tupla:
\[
 F_0(a)=F_0(b)=(8,-2,0,0,-12).
\]
La ruta de prueba tiene
\[
 e(c)=(1,4,1,-2,-4,-3,\;1,4,1,-2,-4,-3)
\]
y
\[
 F_0(c)=(28,-6,0,-4,-12).
\]
Al sumar primero los eventos y reevaluar después, se obtiene
\[
 F_0(a\boxplus c)=(36,-8,0,0,-12),
\]
\[
 F_0(b\boxplus c)=(36,-8,0,-2,-12).
\]
Los dos pares de entradas de la supuesta \(\mu\) son iguales, pero las
salidas difieren en \(\nu_{120}^{\rm ev}\). Esto contradice la existencia de
\(\mu\).
\end{proof}

El censo exhaustivo del mismo dominio contiene \(324\) rutas, \(53\)
cinco-tuplas proxy y \(198\) estados de eventos distintos. Cuarenta y dos
fibras de la proyección contienen más de un estado enriquecido; cuarenta de
ellas son separadas por alguna ruta de prueba. El rombo anterior es el testigo
canónico más breve de esa pérdida de información.

\begin{remark}[Separación tipada entre el representante \(K_0\) y la torsión fuerte]
\label{rem:proxy-no-torsion-fuerte-rev6}
La tercera coordenada de \(F_0\) es \(K_0(e)\), que en B1 toma valores en
\(\{-4,-2,0,2,4\}\). No se identifica con
\(k_{\Delta_4}=T_\zeta-2C_\zeta\). De hecho, existen rutas con el mismo
\(n_A\) y el mismo vector de doce eventos a las que el lector torsional
asigna valores distintos porque conserva los predecesores de cada iteración
y la corona.

La \cref{prop:no-descenso-proyeccion-b1-rev6} demuestra exactamente que la
proyección B1 de cinco coordenadas no transporta la superposición de
registros. No afirma una imposibilidad universal para toda estructura futura
sobre \(\ZZ^5\), ni autoriza a sustituir \(K_0\) por
\(k_{\Delta_4}\). En el lector fuerte, la salida se define mediante la
\cref{def:extractor-posterior-empalme-rev6}: primero se compone la palabra
enriquecida y después se reextraen sus cinco coordenadas.
\end{remark}

\section{Teorema de composición de registros independiente de datos metrológicos}
\label{sec:teorema-composicional-sin-blancos-rev6}

\begin{theorem}[Cierre composicional sin blancos]
\label{thm:cierre-composicional-sin-blancos-rev6}
Fijados el atlas \(\mathcal C_{81}\), la dualidad formal, el residuo de color
y el registro enriquecido del \cref{chap:extension-ruta-caracter}, se cumplen
las afirmaciones siguientes.
\begin{enumerate}
 \item Existe el sistema finito de composiciones tipadas
 \(\mathfrak G_{\rm comp}\) con \(81\) identidades, \(432\) expresiones
 mesónicas, \(1728\) expresiones bariónicas y \(372\) flechas cargadas,
 descompuestas como \(320+52\).

 \item El sistema se construye sin masas objetivo. Su salida es una
 expresión, órbita o multisección tipada, no una masa ni una etiqueta física
 individualizada.

 \item Dada una expresión y dados registros constituyentes con un árbol
 explícito de composiciones compatibles \((g,B)\), la operación produce un
 registro enriquecido. Su firma se obtiene por la aplicación
 \[
 \boxed{
 \bigl(
 \text{expresión},
 \text{registros},
 \text{composiciones orientadas}
 \bigr)
 \xrightarrow{\ \mathfrak C_{12}\ }
 \text{registro compuesto}
 \xrightarrow{\ L_{\rm tick}\ }
 \ZZ^5.}
 \]
 Aquí \(\mathfrak C_{12}\) denota la iteración de
 \(\star_{(g,B)}\) con el parentizado y los bordes incluidos en los datos.

 \item La proyección B1 a cinco enteros no constituye por sí sola un espacio
 composicional: el diagrama con \(\boxplus\) no desciende a
 \(\ZZ^5\). Por tanto, reemplazar cada registro por su cinco-tupla antes del
 composición pierde información necesaria para reevaluar al menos
 \(\nu_{120}^{\rm ev}\).
\end{enumerate}
\end{theorem}

\begin{proof}
El primer punto es el
\cref{thm:censo-gramatica-composicional-rev6}. El segundo es el
\cref{cor:constructor-composicional-sin-blancos}. Para el tercero, cada
arista del árbol de composición se evalúa mediante la
\cref{def:dato-empalme-orientado-rev6}; como los bordes y transportes forman
parte del dominio, la iteración produce un registro sobre el cual está
definido \(L_{\rm tick}\). No se utiliza una firma intermedia como sustituto
del registro. El cuarto punto es la
\cref{prop:no-descenso-proyeccion-b1-rev6}.
\end{proof}

La fórmula correcta de la firma compuesta es, por tanto,
\[
 \operatorname{Sig}_{\rm comp}
 =
 L_{\rm tick}\circ\mathfrak C_{12}.
\]
No es una suma de filas de masa. Sólo después puede aplicarse el carácter
formal o su especialización global:
\[
 \text{registro compuesto}
 \xrightarrow{L_{\rm tick}}\ZZ^5
 \xrightarrow{\chi_{\rm form}}\operatorname{Fun}(\mathcal P,\RR_{>0})
 \longrightarrow
 \text{carta dimensional}.
\]

\section{Reconstrucción genealógica de mezcla, generación y persistencia de sabor}
\label{sec:significado-genealogico-composicion-rev6}

El teorema anterior completa un estrato que una lista de masas no puede
mostrar. Las celdas primitivas no se limitan a recibir nombres; generan
identidades, pares duales, triples de color y flechas orientadas mediante
reglas internas. Los objetos compuestos aparecen así dentro del dominio antes
de toda comparación metrológica.

La genealogía tampoco se reduce al árbol sintáctico. Dos expresiones con el
mismo número de constituyentes pueden componerse de maneras distintas porque
la fase, la hoja, la orientación y la subtraza común forman parte del dato.
El registro conserva esas diferencias hasta que han actuado los signos, las
correlaciones y la torsión. La masa, cuando exista una realización
dimensional, será una lectura posterior de esa memoria compuesta.

Este orden explica por qué \(\ZZ^5\) no es el dominio de composición aunque
sea el dominio del carácter. El grupo \((\ZZ^5,+)\) existe y el carácter es
multiplicativo sobre él; lo que no existe para la proyección B1 es una ley
que recupere desde las cinco-tuplas la superposición que ya olvidó la
colocación de los doce eventos. La linealidad del carácter no devuelve la
información eliminada por el extractor.

\section{Del sistema composicional interno a la realización física}
\label{sec:delimitacion-composicion-especies-rev6}

El \cref{thm:cierre-composicional-sin-blancos-rev6} construye el dominio
anterior a los nombres físicos: expresiones ordenadas, duales, triples de
color y flechas compatibles. Sus cardinales cuentan posibilidades internas
de composición. Una especie física, en cambio, es una representación persistente
con sabor, color, espín, estadística y canal determinados. Confundir ambos
niveles convertiría un censo de expresiones en un falso catálogo de
partículas.

\begin{definition}[Dato de realización física]
\label{def:dato-realizacion-fisica-rev6}
Un dato de realización para una especie \(X\) es la tupla
\[
 \mathfrak D_X=
 \bigl(
 \mathfrak t_X,\mathfrak s_X,
 (\mathcal L_{X_i})_{i=1}^{r},
 \mathcal T_X,(g_X,B_X)
 \bigr),
\]
donde:
\begin{enumerate}
 \item \(\mathfrak t_X\) fija el tipo físico ---elemental, nulo, mesónico,
 bariónico, multiquark o topológico--- y su representación;
 \item \(\mathfrak s_X\) es la multisección o expresión interna realizada;
 \item \(\mathcal L_{X_i}\) son los registros enriquecidos de los
 constituyentes;
 \item \(\mathcal T_X\) fija sabor, color, espín, estadística y
 simetrización;
 \item \((g_X,B_X)\) fija el árbol de composición, sus transportes y sus
 subtrazas comunes.
\end{enumerate}
La tupla se declara sin consultar la masa o la anchura que después se
evaluarán.
\end{definition}

\begin{proposition}[La realización cambia el registro, no la ley]
\label{prop:realizacion-compuesta-ley-unica-rev6}
Todo dato \(\mathfrak D_X\) determina un registro compuesto
\[
 \mathcal L_X
 =\mathfrak C_{12}
   \bigl((\mathcal L_{X_i})_{i=1}^{r};
         \mathcal T_X,g_X,B_X\bigr),
\]
una firma \(v_X=L_{\rm tick}(\mathcal L_X)\) y el carácter central universal
\[
 \chi_0(v_X-v_e).
\]
Si la genealogía aporta además un refinamiento espectral tipado
\(\eta_X\in\mathcal E_{\mathfrak S}\), la evaluación completa es
\[
 m_X
 =m_e^{\rm HMT}\,
   \chi_{\rm full}(v_X-v_e,\eta_X).
\]
El carácter central y el refinamiento pertenecen a niveles distintos; la
completitud de la torre no sustituye el mapa prospectivo que produzca
\(\eta_X\).
Si \(X\) es inestable, su anchura pertenece al lector temporal
\(\Gamma_X=\hbar_{\rm int}/\tau_X^{\rm life}\), no a una sexta coordenada de
la firma de masa.
\end{proposition}

\begin{proof}
El dato de realización contiene exactamente las entradas de
\(\mathfrak C_{12}\). El
\cref{thm:cierre-composicional-sin-blancos-rev6} produce el registro
enriquecido y la \cref{def:extractor-posterior-empalme-rev6} produce su
firma. El carácter central y su elevación espectral son los de
\cref{eq:caracter-masa-global-rev7,eq:ley-masa-todos-registros-rev6}. La
anchura usa una recta temporal y por ello tiene otro dominio.
\end{proof}

La separación de niveles queda así fijada:
\[
\boxed{
\begin{gathered}
\text{sistema composicional interno}
\longrightarrow\text{dato de realización}
\longrightarrow\text{registro compuesto}\\
\longrightarrow\text{firma}
\longrightarrow\text{carácter central}
\longrightarrow\text{refinamiento espectral}
\longrightarrow\text{comparación física}.
\end{gathered}}
\]
Las \(432\) expresiones mesónicas, las \(1728\) bariónicas y las \(372\)
flechas cargadas ocupan sólo la primera casilla y no constituyen un
inventario de partículas. El Modelo Estándar conserva su conjunto tipado de
especies y cada especie ocupa la cadena mediante su propio dato de
realización. La regla causal es única en todos los sectores:
\[
 \boxed{\text{componer extensiones y registros;\quad leer la firma después}.}
\]
